\section{Prediction as a Shared Relation}
\label{sec:prediction}
A shepherd places one stone in a bag for each sheep that enters an enclosure. Before counting the sheep again, the shepherd can use the stones to predict the result. The present stones and the later count refer to the same organized process: a one-to-one registration was made, the registration persisted, and changes to the flock were either prevented or recorded. The example makes prediction concrete. Something available now constrains an outcome that has not yet been read.

Similar organization supports a neural network's forecast, an animal's anticipation, and a machine's coordinated action. Information passes through several forms: a physical relation, a record, a transformation of that record, and a response. Some relations remain distributed among components. Their consequences can be collected through several routes, whose agreement depends on how the components preserve and compare the relevant distinctions. We call this shared organization \emph{synchronicity}.

The statistical foundation is classical information theory. Shannon information measures how records reduce uncertainty about selected alternatives~\citep{shannon1948}. Conditional-expectation geometry resolves which functions of those alternatives are recoverable~\citep{calmon2017}; usable-information theory makes the reader's computational resources explicit~\citep{xu2020}. Our concern is the ordered organization of those recoverable distinctions: how they spread through parallel responses, are compared, and return as a prediction or action.

The geometry in \cref{fig:diamond} is central to this account. An initiating event and a later readout delimit a region of possible propagation. Intermediate responses can coexist across a spatial cut, while each route has an order. Recombination brings those routes back to a common question. The unresolved interior represents continuations that share the available boundary records. Prediction consists of identifying what these continuations already have in common. Phase describes how a distinction changes orientation when carried through their response structure.

\begin{figure*}[t]
\centering
\begin{tikzpicture}[x=1cm,y=1cm,>=Stealth,font=\footnotesize]
\path[use as bounding box] (-.3,-.5) rectangle (16.8,5.35);
\node[anchor=west,font=\bfseries\small] at (0,5.15) {(a) Propagation and recombination};
\fill[routeblue!5] (2.25,.6)--(.35,2.5)--(2.25,4.4)--(4.15,2.5)--cycle;
\draw[gray!55,thin] (2.25,.6)--(1.3,2.5)--(2.25,4.4);
\draw[gray!55,thin] (2.25,.6)--(3.2,2.5)--(2.25,4.4);
\draw[gray!55,thin] (2.25,.6)--(2.25,4.4);
\draw[routeblue,very thick,->] (2.25,.6)--(.35,2.5);
\draw[routeblue,very thick,->] (.35,2.5)--(2.25,4.4);
\draw[routeorange,very thick,dashed,->] (2.25,.6)--(4.15,2.5);
\draw[routeorange,very thick,dashed,->] (4.15,2.5)--(2.25,4.4);
\draw[gray,densely dashed] (.35,2.5)--(4.15,2.5);
\fill (2.25,.6) circle (1.7pt);
\fill (2.25,4.4) circle (1.7pt);
\fill[routeblue] (.35,2.5) circle (1.7pt);
\fill[routeorange] (4.15,2.5) circle (1.7pt);
\node[anchor=north,align=center] at (2.25,.42) {Preparation $p$};
\node[anchor=south,align=center] at (2.25,4.53) {Future readout $q$: target $Z$};
\node[routeblue,anchor=east] at (.2,2.5) {$l$};
\node[routeorange,anchor=west] at (4.3,2.5) {$r$};
\node[routeblue,fill=white,inner sep=1.5pt] at (.95,1.6) {route $\alpha$};
\node[routeorange,fill=white,inner sep=1.5pt] at (3.55,1.6) {route $\beta$};
\node[fill=white,inner sep=2pt,align=center] at (2.25,3.12) {Unresolved\\continuations};
\node[fill=white,inner sep=1pt] at (2.25,2.5) {parallel cut};
\draw[->,thin] (-.16,.6)--(-.16,1.55) node[above] {$t$};
\draw[->,thin] (-.16,.6)--(.65,.6) node[below] {$x$};

\draw[gray!35] (5.05,0)--(5.05,4.95);
\node[anchor=west,font=\bfseries\small] at (5.45,5.15) {(b) Elimination of alternatives};
\node[anchor=west] at (5.55,4.5) {Before the record: $\Omega$};
\foreach \k in {0,...,7}{\draw[fill=gray!15,draw=gray!65] (5.55+.48*\k,3.85) rectangle +(0.35,.35);}
\foreach \k in {0,...,7}{\pgfmathtruncatemacro{\outcome}{\k+1}\node at (5.725+.48*\k,4.025) {\outcome};}
\draw[->,thick] (7.4,3.65)--(7.4,3.15);
\node[anchor=west] at (7.63,3.4) {record $X_c$};
\node[anchor=west] at (5.55,2.95) {Compatible: $\Omega_{c,x}$};
\foreach \k in {0,...,7}{\draw[fill=white,draw=gray!35] (5.55+.48*\k,2.3) rectangle +(0.35,.35);}
\foreach \k in {0,2,3,6}{\draw[fill=routeblue!25,draw=routeblue] (5.55+.48*\k,2.3) rectangle +(0.35,.35);}
\foreach \k in {0,...,7}{\pgfmathtruncatemacro{\outcome}{\k+1}\node[text=gray!55] at (5.725+.48*\k,2.475) {\outcome};}
\foreach \k in {0,2,3,6}{\pgfmathtruncatemacro{\outcome}{\k+1}\node at (5.725+.48*\k,2.475) {\outcome};}
\draw[->,thick] (7.4,2.1)--(7.4,1.6);
\node[anchor=west] at (7.63,1.85) {record $V$};
\foreach \k in {0,...,7}{\draw[fill=white,draw=gray!35] (5.55+.48*\k,.85) rectangle +(0.35,.35);}
\foreach \k in {2,6}{\draw[fill=routeblue!25,draw=routeblue] (5.55+.48*\k,.85) rectangle +(0.35,.35);}
\foreach \k in {0,...,7}{\pgfmathtruncatemacro{\outcome}{\k+1}\node[text=gray!55] at (5.725+.48*\k,1.025) {\outcome};}
\foreach \k in {2,6}{\pgfmathtruncatemacro{\outcome}{\k+1}\node at (5.725+.48*\k,1.025) {\outcome};}
\node[align=center] at (7.45,.25) {Remaining outcomes: $3$ or $7$\\Odd parity is already predicted};

\draw[gray!35] (10.4,0)--(10.4,4.95);
\node[anchor=west,font=\bfseries\small] at (10.85,5.15) {(c) Comparison of predictive routes};
\begin{scope}[yshift=.4cm]
\draw[->,gray] (11.65,1.85)--(15.65,1.85) node[right] {$h_1$};
\draw[->,gray] (13.1,.8)--(13.1,4.35) node[above] {$h_2$};
\draw[gray!35] (13.1,1.85) circle (1.5);
\draw[routeblue,very thick,->] (13.1,1.85)--(14.58,2.11);
\draw[routeorange,very thick,dashed,->] (13.1,1.85)--(13.61,3.26);
\node[routeblue,anchor=west] at (14.6,2.25) {$\U_\alpha h$};
\node[routeorange,anchor=south] at (13.7,3.35) {$\U_\beta h$};
\draw[->,thick] (13.374,2.602) arc[start angle=70,end angle=10,radius=.8];
\node at (13.97,2.58) {$\phi$};
\end{scope}
\node[align=center] at (13.45,.25) {$\Hol_{\beta\leftarrow\alpha}=\U_\beta^*\U_\alpha$\\same question, relative orientation};
\end{tikzpicture}

\caption{Prediction joins distributed evidence to a common future question. \textbf{(a)} A causal diamond joins preparation \(p\) to future readout \(q\); two forward routes \(\alpha,\beta\) sample and combine intermediate responses, while the shaded interior contains unresolved continuations. \textbf{(b)} Records restrict the candidate outcomes to 3 and 7, determining odd parity before the outcome is read. \textbf{(c)} Registering the same predictive distinction along two routes gives a relative transformation \(\Hol_{\beta\leftarrow\alpha}=\U_\beta^*\U_\alpha\); its angle is information phase. The spacetime sketch uses \(c=1\).}
\label{fig:diamond}
\end{figure*}

We develop one mathematical object throughout: a family of conditional-prediction operators together with registration maps between their recoverable modes. Its norm measures predictive strength, its polar factors separate changes of strength from changes of orientation, and comparisons of alternative routes give information phase. The construction makes the meaning of the complex plane explicit: its coordinates are paired future distinctions, and its angle determines which combination a response addresses.

The paper first states the prediction principle and its Shannon specialization (\cref{sec:information}). It then defines complex information and phase within the rhombic geometry, and derives observable prediction-error consequences (\cref{sec:phase,sec:consequences}). Learning appears as the acquisition and reuse of further predictive relations (\cref{sec:learning}). Finally, a spinorial informationon hypothesis gives a common algebraic construction of causal intervals and phase transport (\cref{sec:informationons}). Proofs and calibration procedures are collected in the appendices.


